\chapter*{摘要}
\addcontentsline{toc}{chapter}{摘要}
大模型智能体读取网页、工具返回和长期记忆时，可能把其中的恶意指令当作任务要求，进而读取秘密、外发数据或推送代码。安全研判又需要保留危险原文，单纯删除可疑内容会损失业务证据。

溯界在智能体宿主与工具之间核验动作：以任务契约限定资源与副作用，以来源用途标签约束不可信数据，以放行、审批、阻断三态控制执行，再从事件账本重建溯源图。原型支持真实MCP调用、跨角色委派、跨会话记忆和动作绑定审批；结构化返回保存完整快照，并逐项检查键名与字符串值。凭据目录、Shell启动文件、VCS钩子与IDE设置等敏感目标另受不可覆盖的写入拒绝策略约束。检测器提供风险解释，动作授权由契约、来源与策略共同决定。

提交快照1b64b1f的运行时源码与经CI验证的96d591b一致，Ubuntu、Windows和macOS分别有1107、1106和1107项回归通过。历史48单元离线比较中46项业务完成，两个对照单元实际发生危险推送；规则初筛在140条公开样本上的召回率为56.25\%。这些结果支持原型行为复核，尚不构成真实模型攻击成功率。

\noindent\textbf{关键词：}大模型智能体；提示注入；工具授权；MCP；来源用途约束
\clearpage
